\section{Capacity functionals and active probing}
\label{sec:hit-comparison}
The relation between intersection probabilities and geometry predates
the collision experiment studied here. Matheron~\cite{Matheron} formulates
stereological and morphological information using probabilities that a
structuring figure hits a random set or is contained in it. The capacity
functional $T_X(K)=\Prob\{X\cap K\ne\varnothing\}$ is a central
object in the theory of random closed sets; see Molchanov~\cite{Molchanov},
Chapter~1. The following identity places our observation in that
framework without identifying its endpoint convention with an ordinary
hit test.

\begin{proposition}[The endpoint-excluded capacity datum]
\label{prop:capacity}
Let $Q$ have the commanded density $f$, let $\mathcal O$ be a fixed
closed set, and let $X=\mathcal O-Q$. Then
\begin{equation}\label{eq:capacity-bit}
 P_a(f)=T_X([0,a])-T_X(\{0\}).
\end{equation}
For $X_a=\mathcal O-(Q+a)$ the opposed translated experiment satisfies
\[
 T_{X_a}([-a,0])=T_X([0,a]),\qquad
 P_a(f)-P_{-a}(f_a)=T_{X_a}(\{0\})-T_X(\{0\}).
\]
The equality holds without a convexity or periodicity hypothesis.
\end{proposition}
\begin{proof}
The event $0\in X$ is contained in the event $X\cap[0,a]\ne\varnothing$.
Their difference is precisely a free start with a segment hit. In the
opposed experiment the tested segment has the same physical image, while
its starting endpoint is $Q+a$. Subtraction gives the second formula.
No conditional probability is introduced: solid-start attempts remain
among all draws of $Q$.
\end{proof}

Thus the affine family of capacity probabilities itself is not a new
notion, and the morphology of $C+[-a,0]$ is not a new set operation.
The endpoint exclusion cancels the common segment-hit term and exposes
a signed spatial occupation difference. The bounded zero witness makes
that difference invertible on the specified physical class. The stopped
averaging inverse extends this cancellation to coupled displacement
mixtures. Neither averaging nor a finite-horizon obstacle recursion is
claimed as a new abstract probability principle; the statements specify
which experimentally accessible forcing and which geometric stopping
condition give exact recovery and a finite error bound.

Active probing supplies a second nearby comparison, with different
response alphabets. Edelsbrunner and Skiena~\cite{ES} reconstruct convex
polygons by X-ray probes reporting the \emph{length} of intersection
with a line; their determination bound is $5n+O(1)$ probes for an
$n$-vertex polygon, and their verification bounds concern a given
polygon. A real intersection length is not our Bernoulli collision
output. Their finite-dimensional polygonal target and adaptive exact
line probes are also different from the smooth-class probability
estimation and localized spatial inputs used here. One cannot compare
those probe counts with \eqref{eq:constructive-attempts} while ignoring
both the unknown class and the information returned by a probe.

Bose, De Carufel, Shaikhet and Smid~\cite{BCSS} give a reconstruction
algorithm from wedge probes. Their Theorem~4.1 assumes
$0<\omega\le\pi/2$ and every polygon angle greater than $\omega$,
and obtains at most $2n-2$ probes. A probe reports the wedge position,
its rays and boundary contact points. Our experiment supplies no contact
point and returns one bit per attempt, but pays for a growing grid of
spatially localized commands and repetitions. Consequently neither
observation is asserted to dominate the other in an information order.
These comparisons concern the cited probe models and results, not an
exhaustive priority search in geometric tomography.

The contribution of the present revision is accordingly precise. It
couples the scalar occupation inverse to a stopped randomized-preparation
version, proves a grazing-uniform apparatus-error bound, and replaces
one infinite-class existence selection by a proved finite geometric
procedure. The bounded-period and patch-margin hypotheses remain visible
in the global conclusion. No statistical lower bound or optimality,
passive rigidity theorem, discovery of periodicity without a prior, or
exact analytic uniform-count theorem is inferred from these additions.

\subsection*{Current submission map}
The five core chapters of the reviewed scalar article remain active and
unchanged. The complete reviewed article, including its validation files
and all nested earlier volumes, is also supplied at \texttt{retained/v29}.
In that retained terminology, Supplement P denotes v28; its current path
is \texttt{retained/v29/retained/v28}. The present primary adds the
calibration, stopped inverse, constructive recovery and capacity comparison
as self-contained arguments. Earlier existence-only statements describe
their original algorithms and are not relabelled as computational bounds
for this revision. The source map is a preservation record, not a claim
that a former build receipt validates the new proofs.
